\begin{helper}
Let \(s\ge4\), let \(m\ge1\), put \(t=s-2\) and \(q=2^m\).  Define
\[
  n=\frac{q^{t+1}-1}{q-1},\qquad
  d_G=\frac{q^t-1}{q-1},\qquad
  d_F=n-d_G,\qquad
  a=\frac{q^{t-1}-1}{q-1},
\]
and
\[
  \lambda=\sqrt{d_G-a},\qquad Q=q^{t-1}.
\]
Then \(n,d_F,d_G\) are positive,
\[
  q^t\le n\le (t+1)q^t,\qquad
  q^{t-1}\le d_G\le t q^{t-1},
\]
\[
  d_F=q^t,\qquad d_G-a=q^{t-1},\qquad \lambda^2=Q,
\]
and \(Q>0\).  Moreover
\[
  \frac1{n^2}\le \frac{\lambda^2}{d_G^2},\qquad
  \frac{\lambda^2}{d_G^2}\le 1,\qquad
  \frac{\lambda\lambda}{d_Fd_G}\le 1,
\]
and
\[
  \max\left\{\frac{\lambda^2}{d_G^2},
    \frac{\lambda\lambda}{d_Fd_G}\right\}\le \frac1Q.
\]
\end{helper}
\begin{proof}
Since \(s\ge4\), we have \(t=s-2\ge2\).  Since \(m\ge1\), the number
\(q=2^m\) satisfies \(q\ge2\).

For any \(r\ge1\), the finite geometric-series identity gives
\[
  1+q+\cdots+q^{r-1}=\frac{q^r-1}{q-1}.
\]
Applying this with \(r=t+1,t,t-1\), respectively, gives
\[
  n=\sum_{i=0}^{t}q^i,\qquad
  d_G=\sum_{i=0}^{t-1}q^i,\qquad
  a=\sum_{i=0}^{t-2}q^i.
\]
Therefore
\[
  n=d_G+q^t,\qquad d_G=a+q^{t-1}.
\]
It follows immediately that
\[
  d_F=n-d_G=q^t,\qquad d_G-a=q^{t-1}.
\]
The displayed formulas also give \(n>0\), \(d_F>0\), and \(d_G>0\).

The lower bounds \(q^t\le n\) and \(q^{t-1}\le d_G\) follow from
\(n=d_G+q^t\) and \(d_G=a+q^{t-1}\).  For the upper bounds, every term
\(q^i\) with \(0\le i\le t\) is at most \(q^t\), because \(q\ge1\).  Summing
over the \(t+1\) terms in \(n\) gives \(n\le(t+1)q^t\).  Similarly, every
term \(q^i\) with \(0\le i\le t-1\) is at most \(q^{t-1}\), and there are
\(t\) such terms, so \(d_G\le t q^{t-1}\).

Since \(d_G-a=q^{t-1}\), the definition of \(\lambda\) gives
\[
  \lambda^2=(\sqrt{q^{t-1}})^2=q^{t-1}=Q.
\]
Also \(Q>0\), since \(q>0\).

It remains to check the real-valued ratio estimates.  From \(d_G\le n\) and
positivity, \(d_G^2\le n^2\), hence \(1/n^2\le1/d_G^2\).  Since \(Q\ge1\),
\[
  \frac1{d_G^2}\le \frac{Q}{d_G^2}
    =\frac{\lambda^2}{d_G^2},
\]
which proves the first ratio lower bound.

The inequality \(Q\le d_G\) and \(d_G\ge1\) imply \(Q\le d_G\le d_G^2\).
Thus
\[
  \frac{\lambda^2}{d_G^2}=\frac{Q}{d_G^2}\le1.
\]
Also \(d_F\ge1\), so \(d_G\le d_Fd_G\).  Since \(Q\le d_G\),
\[
  \frac{\lambda\lambda}{d_Fd_G}
    =\frac{\lambda^2}{d_Fd_G}
    =\frac{Q}{d_Fd_G}\le1.
\]

For the sharper upper bound, \(Q\le d_G\) gives \(Q^2\le d_G^2\), and hence
\[
  \frac{\lambda^2}{d_G^2}
    =\frac{Q}{d_G^2}\le\frac1Q.
\]
Furthermore \(d_F=q^t\ge q^{t-1}=Q\) and \(d_G\ge Q\), so
\[
  Q^2\le d_Fd_G.
\]
Consequently
\[
  \frac{\lambda\lambda}{d_Fd_G}
    =\frac{Q}{d_Fd_G}\le\frac1Q.
\]
Taking the maximum of the two estimates proves the final displayed
inequality.
\end{proof}
